\section{\titleTightness}\label{sec:tightness}
  \hhnote{ 
    Either insert after O2H section as an ``Interlude'', or leave as an appendix. 
  }
  \hhnote{
    Note: The brute-force setting will still be relevant in a world
    with cryptographically-relevant quantum computers, given how much cheaper
    classical attack resources are. (And also, for many of the relevant
    problems, the best classical attacks see little speed-up from Grover.)
  }

  % Usage: \crackzfigcode{ name of file in ./crakz/ -- gets label "fig:[file name]" }{ caption }
  %\crakzfigcode{qh-attacks}{
  %  The effect of a loss factor $q \leq 2^{80}$ on the 
  %  attacks our security proof rules out, under the assumption that the
  %  underlying primitive achieves $128$-bit security, and that the best attack
  %  on the primitive is a classical brute-force search (left) and Grover
  %  search (right), as a function of the upper-bound $t$ on the adversary's runtime. 
  %  Higher up on the $y$-axis $=$ higher probability $\epsilon$ of a successful break.
  %}

  \subsection{Kinds of Loss: Additive vs.\ Multiplicative vs.\ Roots}
    Say you have a construction built on top of a primitive $\Pi$, that incurs
      some tightness loss. For illustration, we study how this tightness loss
      affects security under the assumption that the best-known attack against
      the primitive is either brute-force search (linear, so $\epsilon_\Pi =
      t/{2^{-\lambda}}$), Grover search or (classical) collision-finding
      (quadratic, $\epsilon_\Pi = t^2/2^{-\lambda}$), or quantum
      collision-finding (cubic, $\epsilon_\Pi = t^3/2^{-\lambda}$). 

    \crefrange{fig:qh-loss}{fig:sq-loss} show how different kinds of tightness
      losses degrade the provable security of our construction, under different
      assumptions. Here, the primitive $\Pi$ is assumed to achieve
      $128$-bit security against various assumed-optimal attacks: 
      Brute-force search in \cref{fig:qh-loss},
      Grover search and/or classical collision-finding in \cref{fig:sq-loss},
      and quantum collision-finding in \cref{fig:cube-loss}.
      \cref{tab:compensation} shows how high one would have to set the security
      parameter $\lambda$ of the construction in order to sufficiently
      compensate for the tightness so that the construction achieves $128$-bit
      security.

    Note that this is not a particularly realistic setting, as even 
      primitives that are considered highly secure (such as AES) 
      usually face attacks that are
      significantly better than the generic ones we consider here.
      Further, different attacks strategies may be optimal for different ranges
      of the time budget $t$. 

    %\crakzfigcode{add-loss}{
    \hhnote{OLD CAPTION 1:
      The effect of an additive loss $q/2^{n} \leq 2^{80-128} = 2^{-48}$.
    }

    %\crakzfigcode{qh-loss}{
    \hhnote{OLD CAPTION 2:
      Attack advantages ruled out by a reduction from the security of a primitive
      $\Pi$ to the security of a construction, if the reduction incurred a $q \leq 2^{80}$
      loss factor in the advantage (and no loss in the running time) (top), and
      the bit-security implied for the construction (bottom), assuming that the
      best attack on the primitive is scales linearly as in a brute-force search
      (left) and quadratically as with Grover search and collision-finding
      (right). We assume that the primitive achieves $128$-bit security against
      the relevant attack, and display the ruled-out advantages/bit security as a
      function of the bound on the adversarial run-time $t$, to highlight the
      time-probability trade-off that adversaries generally face. 
      (So for example, if high $t$ gives lower bit-security, the adversary will be
      incentivized to use all the time available to it, while if the lowest
      bit-security sits at low $t$, the optimal attack may involve performing the
      same low-advantage attack a large number of times in parallel.)
    }

    %\crakzfigcode{sq-loss}{
    \hhnote{OLD CAPTION 3:
      Same as \cref{fig:qh-loss}, except if the reduction incurred a
      square-root loss instead of a factor $q$.
      Interestingly, if the best attack on the primitive already scales
      quadratically, then a square-root yields the same bit-security, the only
      difference being that now we can't rule out matching attacks for any values
      of $t$ (where previously the only matching attack was at $t = t_{\text{max}}
      = 2^{128}$). Contrast with the linear-scaling case, where a square-root
      loss would halve the bit-security.
    }

    %\crakzfigcode{loss-attacks}{
    %  ATTACKS: testing a different combination: loss types left--right, attack types
    %  up--down
    %}
    %\crakzfigcode{loss-bitsec}{
    %  BIT-SECURITY: testing a different combination: loss types left--right, attack types
    %  up--down
    %}

    % try with every loss = 64 bits
    \crakzfigcode{sixtyfourloss-attacks}{
      The $t$-vs.-$\epsilon$ attack landscape for a construction built from some
      primitive $\Pi$, if the security reduction incurred 1) the
      additive loss $+c=2^{-64}$ (top row), 2) the multiplicative loss
      $\cdot\, c=2^{64}$ (middle row), or 3) a square-root loss, 
      assuming that the optimal advantage against $\Pi$ scales a)
      linearly with time, e.g.\ brute-force search (left column), b)
      quadratically with time, e.g.\ Grover
      search/classical collision-finding (middle column), and 
      c) cubically with time, e.g.\ quantum collision-finding (right column).
      We assume that $\Pi$ achieves $128$-bit security against
      the relevant attack type, and that the reduction did not incur any significant
      increase in runtime.
      Dots mark the corresponding bit-security, see \cref{fig:sixtyfourloss-bitsec}.
    }
    \crakzfigcode{sixtyfourloss-bitsec}{
      The bit-security of the primitive $\Pi$ for each kind of assumed-best
      attack, and the implied bit-security of the construction for a given reduction type, as
      a function of the upper-bound $t$ on the run-time of the adversary. 
      The panels correspond to the panels of \cref{fig:sixtyfourloss-attacks},
      and dots mark the assumed/implied bit-security when $t \leq 2^{128}$.
      Observe how a square-root loss does not lead to a loss in bit-security
      once the underlying assumed-optimal attack scales quadratically or better.
    }

  \subsection{Understanding the O2H Losses}
    See \cref{tab:compensation} and \cref{fig:sqO2H-loss}.

    \begin{table}[h]
      \centering
      %
      \caption{
        \label{tab:compensation}
        Compensating for tightness losses: Assuming the primitive $\Pi$ achieves
        $128$-bit security against Attack (and assuming that Attack
        is the most effective attack against the primitive), how large should we set
        the security parameter of the construction to compensate for the security
        loss and provably achieve $128$-bit security?
        Here, as elsewhere, we assume $q \leq 2^{80}$.
        Note how once we assume a primitive attack that scales like $\epsilon
        \propto t^c$ for $c \geq 2$, then a square-root loss no longer needs to be
        compensated for.
      }
      %
      \rowcolors{2}{}{gray!10}
      \begin{tabular}{lc||l|l|l}
        \toprule
        \multicolumn{5}{c}{%
          \textbf{For $128$-bit security against Attack on $\Pi$, set $\lambda$ to \ldots}
        } \\
        \midrule
          & \textbf{Tightness loss}
            & Brute-force search 
              & Grover search/collision finding 
                & Quantum collision finding \\
        \textbf{O2H Variant}
          & $\epsilon \lesssim \ldots$ 
            & $\quad \To \epsilon_\Pi \propto t$, $\lambda_\Pi = 128$ 
              & $\quad \To \epsilon_\Pi \propto t^2$, $\lambda_\Pi = 256$  
                & $\quad \To \epsilon_\Pi \propto t^3$, $\lambda_\Pi = 384$ \\
        \midrule
        \midrule
        Classical Find
          & $\epsilon_\Pi$ 
            & $\lambda \assign 128\ (= \lambda_\Pi)$ 
              & $\lambda \assign 256\ (= \lambda_\Pi)$ 
                & $\lambda \assign 384\ (= \lambda_\Pi)$ \\
        MRE-Find
          & $\sqrt{q} \cdot \epsilon_\Pi$ 
            & $\lambda \assign 168\ (= \lambda_\Pi + 40)$ 
              & $\lambda \assign 296\ (= \lambda_\Pi + 40)$ 
                & $\lambda \assign 424\ (= \lambda_\Pi + 40)$ \\
        Classical Guess
          & $q \cdot \epsilon_\Pi$ 
            & $\lambda \assign 208\ (= \lambda_\Pi + 80)$ 
              & $\lambda \assign 336\ (= \lambda_\Pi + 80)$ 
                & $\lambda \assign 464\ (= \lambda_\Pi + 80)$ \\
        \midrule
        DS-Find
          & $\sqrt{\epsilon_\Pi}$ 
            & $\lambda \assign 256\ (= 2 \cdot \lambda_\Pi)$ 
              & $\lambda \assign 256\ (= \lambda_\Pi)$ 
                & $\lambda \assign 384\ (= \lambda_\Pi)$ \\
        SC-Find
          & $\sqrt{q \cdot \epsilon_\Pi}$ 
            & $\lambda \assign 336\ (= 2 \cdot \lambda_\Pi + 80)$ 
              & $\lambda \assign 336\ (= \lambda_\Pi + 80)$ 
                & $\lambda \assign 464\ (= \lambda_\Pi + 80)$ \\
        Guess
          & $q\sqrt{\epsilon_\Pi}$ 
            & $\lambda \assign 416\ (= 2 \cdot (\lambda_\Pi + 80))$ 
              & $\lambda \assign 416\ (= \lambda_\Pi + 2 \cdot 80)$ 
                & $\lambda \assign 544\ (= \lambda_\Pi + 2 \cdot 80)$ \\
        \bottomrule
      \end{tabular}
    \end{table}

    %\crakzfigcode{GuessO2H-loss}{
    %  $\epsilon \sim q \sqrt{\epsilon_\Pi}$
    %}

    %\crakzfigcode{FindO2H-loss}{
    %  Same as \cref{fig:qh-loss} -- \cref{fig:GuessO2H-loss}, except if the reduction
    %  incurred the Find-O2H loss $\epsilon \sim \sqrt{q \cdot \epsilon_\Pi}$.
    %}

    %\crakzfigcode{O2H-loss}{
    %  Find and Guess combined
    %}

    \crakzfigcode{sqO2H-loss}{
      Same as \cref{fig:qh-loss} and \cref{fig:sq-loss}, except if the reduction
      incurred each of the O2H losses 
      $\epsilon \sim \sqrt{\epsilon_\Pi}$ (line),  
      $\epsilon \sim \sqrt{q \cdot \epsilon_\Pi}$ (line),  
      $\epsilon \sim q \sqrt{\epsilon_\Pi}$ (line).
      Observe how, if parameter were chosen to resist a linearly-scaling attack with
      $128$-bit security, then the tightness loss is large enough to completely void the
      reduction bound, meaning the theorem would imply no security for the
      construction in this case.
    }

